\documentclass{amsart}
% For dealing with references we use the comment environment
\usepackage{verbatim}
\newenvironment{reference}{\comment}{\endcomment}
%\newenvironment{reference}{}{}
\newenvironment{slogan}{\comment}{\endcomment}
\newenvironment{history}{\comment}{\endcomment}

% For commutative diagrams we use Xy-pic
\usepackage[all]{xy}

% We use 2cell for 2-commutative diagrams.
\xyoption{2cell}
\UseAllTwocells

% We use multicol for the list of chapters between chapters
\usepackage{multicol}

% This is generally recommended for better output
\usepackage{lmodern}
\usepackage[T1]{fontenc}

% For cross-file-references

% Package for hypertext links:
\usepackage{hyperref}

% For any local file, say "hello.tex" you want to link to please

% Theorem environments.
%
\theoremstyle{plain}
\newtheorem{theorem}[subsection]{Theorem}
\newtheorem{proposition}[subsection]{Proposition}
\newtheorem{lemma}[subsection]{Lemma}

\theoremstyle{definition}
\newtheorem{definition}[subsection]{Definition}
\newtheorem{example}[subsection]{Example}
\newtheorem{exercise}[subsection]{Exercise}
\newtheorem{situation}[subsection]{Situation}

\theoremstyle{remark}
\newtheorem{remark}[subsection]{Remark}
\newtheorem{remarks}[subsection]{Remarks}

\numberwithin{equation}{subsection}

% Macros
%
\def\lim{\mathop{\mathrm{lim}}\nolimits}
\def\colim{\mathop{\mathrm{colim}}\nolimits}
\def\Spec{\mathop{\mathrm{Spec}}}
\def\Hom{\mathop{\mathrm{Hom}}\nolimits}
\def\Ext{\mathop{\mathrm{Ext}}\nolimits}
\def\SheafHom{\mathop{\mathcal{H}\!\mathit{om}}\nolimits}
\def\SheafExt{\mathop{\mathcal{E}\!\mathit{xt}}\nolimits}
\def\Sch{\mathit{Sch}}
\def\Mor{\mathop{\mathrm{Mor}}\nolimits}
\def\Ob{\mathop{\mathrm{Ob}}\nolimits}
\def\Sh{\mathop{\mathit{Sh}}\nolimits}
\def\NL{\mathop{N\!L}\nolimits}
\def\CH{\mathop{\mathrm{CH}}\nolimits}
\def\proetale{{pro\text{-}\acute{e}tale}}
\def\etale{{\acute{e}tale}}
\def\QCoh{\mathit{QCoh}}
\def\Ker{\mathop{\mathrm{Ker}}}
\def\Im{\mathop{\mathrm{Im}}}
\def\Coker{\mathop{\mathrm{Coker}}}
\def\Coim{\mathop{\mathrm{Coim}}}

% Boxtimes
%
\DeclareMathSymbol{\boxtimes}{\mathbin}{AMSa}{"02}

%
% Macros for moduli stacks/spaces
%
\def\QCohstack{\mathcal{QC}\!\mathit{oh}}
\def\Cohstack{\mathcal{C}\!\mathit{oh}}
\def\Spacesstack{\mathcal{S}\!\mathit{paces}}
\def\Quotfunctor{\mathrm{Quot}}
\def\Hilbfunctor{\mathrm{Hilb}}
\def\Curvesstack{\mathcal{C}\!\mathit{urves}}
\def\Polarizedstack{\mathcal{P}\!\mathit{olarized}}
\def\Complexesstack{\mathcal{C}\!\mathit{omplexes}}
% \Pic is the operator that assigns to X its picard group, usage \Pic(X)
% \Picardstack_{X/B} denotes the Picard stack of X over B
% \Picardfunctor_{X/B} denotes the Picard functor of X over B
\def\Pic{\mathop{\mathrm{Pic}}\nolimits}
\def\Picardstack{\mathcal{P}\!\mathit{ic}}
\def\Picardfunctor{\mathrm{Pic}}
\def\Deformationcategory{\mathcal{D}\!\mathit{ef}}

\setlength{\emergencystretch}{3em}
\begin{document}
\title[Stacks Project, Tag 07M7]{307 --- Etale-local lifting of sections of smooth algebras}
\maketitle
\noindent Copyright (C) 2005--2025 Johan de Jong. Stacks Project authors. Source: \href{https://stacks.math.columbia.edu/tag/07M7}{Tag 07M7}.

\noindent Editorial difficulty note (not author text). Difficulty H2: The proof stabilizes a projective conormal module to a free one by adjoining a symmetric algebra. It chooses equations, verifies vanishing of the cotangent obstruction and differentials, and localizes the resulting algebra to construct an etale lift with the prescribed closed-fibre section..

\noindent Editorial provenance note (not author text). History: excerpt from revision \texttt{a0444\allowbreak 6e57e\allowbreak c1fbc\allowbreak 252a8\allowbreak 71afc\allowbreak ec775\allowbreak 2fb28\allowbreak 07b14}, more-algebra.tex, lines 2138--2227. Reference presentation was adapted; original-author context and core support blocks were retained or added. The mathematical wording is unchanged. Licensed under GNU FDL 1.2 or later; no invariant sections or cover texts. See ../licenses/COPYING. Context blocks reproduce author definitions and supporting results; they are not additional counted problems.

\begin{lemma}
\label{lemma-lift-section-smooth-morphism}
Let $A$ be a ring, let $I \subset A$ be an ideal. Consider a commutative
diagram
$$
\xymatrix{
B \ar[rd] \\
A \ar[u] \ar[r] & A/I
}
$$
where $B$ is a smooth $A$-algebra. Then there exists an \'etale ring
map $A \to A'$ which induces an isomorphism $A/I \to A'/IA'$ and an
$A$-algebra map $B \to A'$ lifting the ring map $B \to A/I$.
\end{lemma}

\begin{proof}
Let $J \subset B$ be the kernel of $B \to A/I$ so that $B/J = A/I$. By
Algebra, Lemma \href{https://stacks.math.columbia.edu/tag/06A9}{lemma application NL smooth (Tag 06A9)} the sequence
$$
0 \to I/I^2 \to J/J^2 \to \Omega_{B/A} \otimes_B B/J \to 0
$$
is split exact. Thus $\overline{P} = J/(J^2 + IB) = \Omega_{B/A} \otimes_B B/J$
is a finite projective $A/I$-module. Choose an integer $n$ and a direct sum
decomposition $A/I^{\oplus n} = \overline{P} \oplus \overline{K}$.
By Lemma \href{https://stacks.math.columbia.edu/tag/07M5}{lemma lift projective module (Tag 07M5)} we can find an
\'etale ring map $A \to A'$ which induces an isomorphism
$A/I \to A'/IA'$ and a finite projective $A$-module $K$ which
lifts $\overline{K}$. We may and do replace $A$ by $A'$.
Set $B' = B \otimes_A \text{Sym}_A^*(K)$. Since $A \to \text{Sym}_A^*(K)$
is smooth by Lemma \href{https://stacks.math.columbia.edu/tag/07M6}{lemma symmetric algebra smooth (Tag 07M6)} we see that
$B \to B'$ is smooth which in turn implies that $A \to B'$ is smooth (see
Algebra, Lemmas \href{https://stacks.math.columbia.edu/tag/00T4}{lemma base change smooth (Tag 00T4)} and
\href{https://stacks.math.columbia.edu/tag/00TC}{lemma locally smooth (Tag 00TC)}).
Moreover the section $\text{Sym}^*_A(K) \to A$ determines a section
$B' \to B$ and we let $B' \to A/I$ be the composition $B' \to B \to A/I$.
Let $J' \subset B'$ be the kernel of $B' \to A/I$. We have
$JB' \subset J'$ and $B \otimes_A K \subset J'$. These maps combine
to give an isomorphism
$$
(A/I)^{\oplus n} \cong J/(J^2 + IB) \oplus \overline{K}
\longrightarrow
J'/((J')^2 + IB')
$$
Thus, after replacing $B$ by $B'$ we may assume that
$J/(J^2 + IB) = \Omega_{B/A} \otimes_B B/J$ is a free
$A/I$-module of rank $n$.

\medskip\noindent
In this case, choose $f_1, \ldots, f_n \in J$ which map to a
basis of $J/(J^2 + IB)$. Consider the finitely presented $A$-algebra
$C = B/(f_1, \ldots, f_n)$. Note that we have an exact sequence
$$
0 \to H_1(L_{C/A}) \to (f_1, \ldots, f_n)/(f_1, \ldots, f_n)^2
\to \Omega_{B/A} \otimes_B C \to \Omega_{C/A} \to 0
$$
see Algebra, Lemma \href{https://stacks.math.columbia.edu/tag/00S2}{lemma exact sequence NL (Tag 00S2)} (note that
$H_1(L_{B/A}) = 0$ and that $\Omega_{B/A}$ is finite projective,
in particular flat so the Tor group vanishes). For any prime
$\mathfrak q \supset J$ of $B$ the module $\Omega_{B/A, \mathfrak q}$
is free of rank $n$ because $\Omega_{B/A}$ is finite projective and
because $\Omega_{B/A} \otimes_B B/J$ is free of rank $n$
(see Algebra, Lemma \href{https://stacks.math.columbia.edu/tag/00NX}{lemma finite projective (Tag 00NX)}). By our choice
of $f_1, \ldots, f_n$ the map
$$
\left((f_1, \ldots, f_n)/(f_1, \ldots, f_n)^2\right)_{\mathfrak q}
\to
\Omega_{B/A, \mathfrak q}
$$
is surjective modulo $J$. Hence we see that this map of modules over
the local ring $C_{\mathfrak q}$ has to be an isomorphism
(by Algebra, Lemma \href{https://stacks.math.columbia.edu/tag/00DV}{lemma NAK (Tag 00DV)} the map is surjective
and for example by Algebra, Lemma \href{https://stacks.math.columbia.edu/tag/05G8}{lemma fun (Tag 05G8)}
because $((f_1, \ldots, f_n)/(f_1, \ldots, f_n)^2)_{\mathfrak q}$
is generated by $n$ elements the map is injective). Thus
$H_1(L_{C/A})_{\mathfrak q} = 0$ and $\Omega_{C/A, \mathfrak q} = 0$. By
Algebra, Lemma \href{https://stacks.math.columbia.edu/tag/07BU}{lemma smooth at point (Tag 07BU)}
we see that $A \to C$ is smooth at the prime $\overline{\mathfrak q}$
of $C$ corresponding to $\mathfrak q$. Since
$\Omega_{C/A, \mathfrak q} = 0$ it is actually \'etale at
$\overline{\mathfrak q}$. Thus $A \to C$ is \'etale at all primes
of $C$ containing $JC$. By Lemma \href{https://stacks.math.columbia.edu/tag/07M0}{lemma localize upstairs (Tag 07M0)}
we can find an $f \in C$ mapping to an invertible element of $C/JC$ such that
$A \to C_f$ is \'etale. By our choice of $f$ it is still true that
$C_f/JC_f = A/I$. The map $C_f/IC_f \to A/I$ is surjective and
\'etale by Algebra, Lemma \href{https://stacks.math.columbia.edu/tag/00U7}{lemma map between etale (Tag 00U7)}.
Hence $A/I$ is isomorphic to the localization of $C_f/IC_f$ at
some element $g \in C$, see
Algebra, Lemma \href{https://stacks.math.columbia.edu/tag/00U8}{lemma surjective flat finitely presented (Tag 00U8)}.
Set $A' = C_{fg}$ to conclude the proof.
\end{proof}
\end{document}
